% Introduction
\section{Introduction}

Conversational AI systems are trained to be helpful, but what makes a conversation genuinely helpful remains poorly understood. While current alignment methods focus on training AI to match human preferences through techniques like Reinforcement Learning from Human Feedback (RLHF) \cite{ouyang2022training, bai2022constitutional}, they treat alignment as a unidirectional optimization problem: adjust the AI policy until humans rate it favorably. This framing ignores a fundamental aspect of successful conversations observed in human-human dialogue---both participants adapt to each other over time \cite{pickering2004toward}. When users interact with AI assistants across multiple turns, they are not static evaluators but active learners who refine their communication strategies based on what works.

Consider what happens during a typical multi-turn conversation with an AI assistant. A user begins with an initial query, receives a response, and then decides how to proceed. If the response was unhelpful, they reformulate the question---rephrasing, adding context, or breaking it into sub-questions. Over successive turns, experienced users learn which query types the AI handles well, how to frame requests effectively, and where the AI's capabilities end. Simultaneously, the AI system---if designed with responsiveness mechanisms---may adapt its output style, detail level, or exploration depth based on the user's query patterns. This bidirectional co-adaptation, where both user and AI adjust their behaviors during interaction, suggests that alignment is not a static property of the AI policy alone but an emergent property of the conversation itself.

Current alignment evaluation methods fail to capture this interactive dynamic. RLHF optimizes for high human ratings by adjusting model parameters during training, but once deployed, the model's policy is fixed. Evaluation focuses on single-turn quality or aggregate satisfaction scores, not on how users and AI co-adapt within conversations. We lack metrics that distinguish between conversations where only the AI is ``aligned'' (produces acceptable responses) and conversations where genuine bidirectional alignment occurs---users learn to communicate effectively AND the AI exhibits responsiveness to evolving query patterns.

We propose \textbf{behavioral coupling} as a metric for bidirectional alignment in conversational AI. The key insight is that co-adaptation leaves observable behavioral signatures: users who learn AI capabilities should reformulate queries less over turns (learning curve effect), while responsive AI systems should produce outputs whose diversity correlates with query diversity. By measuring the correlation between these two signals---user learning rate (reformulation slope) and AI responsiveness (query-response diversity correlation)---we quantify the strength of bidirectional alignment in individual conversations.

To validate this framework, we analyze 169,352 multi-turn conversations from the HH-RLHF dataset \cite{bai2022constitutional}, focusing on conversations with $\geq 5$ turns where learning curves can be detected. We test two existence hypotheses: (1) users exhibit decreasing reformulation rates over turns, indicating learning of AI capabilities, and (2) AI response diversity correlates with query diversity, indicating responsiveness to user patterns. Our results show:

\begin{itemize}
\item \textbf{User learning exists}: Reformulation rates decrease over conversation turns with mean slope = $-0.021$ ($p=0.012$, Cohen's $d=-0.246$), confirming users adapt their query strategies during interactions.
\item \textbf{AI responsiveness exists}: Response diversity correlates with query diversity ($r=0.396$, $p<0.001$, 95\% CI [0.392, 0.400]), indicating the AI system tracks user query patterns.
\item \textbf{Temporal co-adaptation}: The diversity correlation strengthens with conversation length ($r=0.04$ for 2--3 turns $\rightarrow$ $r=0.19$ for 6+ turns), consistent with alignment emerging over extended interactions rather than being present from the first turn.
\end{itemize}

While effect sizes are smaller than initially hypothesized (small Cohen's $d$ for learning, correlation below our $r>0.4$ target threshold), both components of bidirectional alignment are statistically robust across a large sample. These findings validate the theoretical framework and provide a foundation for future work testing whether coupling strength---the interaction between user learning and AI responsiveness---predicts conversation quality beyond individual metrics.

The contribution of this work is threefold: \textbf{(1)} We introduce behavioral coupling as an operationalization of bidirectional alignment, shifting measurement from static policy evaluation to conversation dynamics. \textbf{(2)} We provide empirical evidence that both components of coupling (user learning and AI responsiveness) exist in RLHF-trained systems, validating the framework's theoretical foundation. \textbf{(3)} We demonstrate that these behavioral signals can be extracted from standard conversation logs without additional human evaluation, enabling scalable alignment monitoring in deployed systems.

The remainder of the paper is organized as follows: Section~\ref{sec:related} reviews related work in AI alignment methods, user adaptation in HCI, and conversation-level quality metrics. Section~\ref{sec:methodology} describes our methodology for detecting reformulation patterns and measuring diversity correlations. Section~\ref{sec:experiments} presents experimental design and datasets. Section~\ref{sec:results} reports results from our validation experiments. Section~\ref{sec:discussion} discusses implications, limitations, and future directions. Section~\ref{sec:conclusion} concludes with a vision for bidirectional alignment training methods.
